\chapter{Elaborazione \& Analisi}

\section{Elaborazione dei dati}

I dati ottenuti durante la fase di testing, salvati nei file \texttt{hash\_data.csv} e \texttt{op\_data.csv} per il test con potenze di 10 e i file 
\texttt{hash\_data\_double.csv} e \texttt{op\_data\_double.csv} per il test al raddoppio, sono stati successivamente rielaborati mediante l'utilizzo 
del linguaggio Python. In particolare, sono state impiegate le librerie \texttt{pandas} per la manipolazione e l'analisi dei dati e \texttt{matplotlib.pyplot} per la generazione dei grafici.

Sono stati sviluppati quattro script principali:
\begin{itemize}
    \item \texttt{grapher\_hash.py} e \texttt{grapher\_op.py}, dedicati all'analisi dei dati relativi a ciascuna singola tipologia;
    \item \texttt{grapher\_mix\_hash.py} e \texttt{grapher\_mix\_op.py}, utilizzati per confrontare i risultati tra diverse tipologie di dati.
\end{itemize}

Per ciascuna tipologia di dati sono stati generati:
\begin{itemize}
    \item otto grafici relativi alle configurazioni della hashtable;
    \item nove grafici dedicati all'analisi delle singole operazioni.
\end{itemize}

In aggiunta, sono stati realizzati cinque grafici complessivi che mettono a confronto le prestazioni tra le diverse tipologie di dati analizzate.

Gli otto grafici relativi alla configurazione della hashtable sono:
\begin{itemize}
    \item \texttt{latency\_by\_hash\_strategy.pdf}
    \item \texttt{latency\_by\_max\_load\_factor.pdf}
    \item \texttt{latency\_by\_overflow\_strategy.pdf}
    \item \texttt{latency\_by\_resize\_strategy.pdf}
    \item \texttt{memory\_by\_hash\_strategy.pdf}
    \item \texttt{memory\_by\_max\_load\_factor.pdf}
    \item \texttt{memory\_by\_overflow\_strategy.pdf}
    \item \texttt{memory\_by\_resize\_strategy.pdf}
\end{itemize}

I nove grafici relativi alle operazioni sono:
\begin{itemize}
    \item \texttt{key\_distribution\_by\_overflow.pdf}
    \item \texttt{latency\_by\_operations\_hash\_strategy.pdf}
    \item \texttt{latency\_by\_operations\_\&\_max\_load\_factor.pdf}
    \item \texttt{latency\_by\_operations\_\&\_overflow\_strategy.pdf}
    \item \texttt{latency\_by\_operations\_\&\_resize\_strategy.pdf}
    \item \texttt{latency\_insert\_exist\_by\_load\_factor.pdf}
    \item \texttt{memory\_by\_operations\_\&\_max\_load\_factor.pdf}
    \item \texttt{memory\_by\_operations\_\&\_overflow\_strategy.pdf}
    \item \texttt{memory\_by\_operations\_\&\_resize\_strategy.pdf}
\end{itemize}

I cinque grafici relativi alle tipologie sono:
\begin{itemize}
    \item \texttt{latency\_by\_keytype\_hash.pdf}
    \item \texttt{latency\_by\_operations\_\&\_max\_load\_factor.pdf}
    \item \texttt{latency\_by\_keytype\_\_overflow.pdf}
    \item \texttt{latency\_by\_keytype\_resize.pdf}
    \item \texttt{memory\_by\_keytype.pdf}
    \item \texttt{latency\_by\_operation\_keytype\_elements.pdf}
\end{itemize}

\begin{lstlisting}[language=Python, caption=Esempio di script per la generazione di grafici]
# -----------------------------
# Latenza / Overflow
# -----------------------------
df = pd.read_csv(path)

# -----------------------------
# Calcola media, minimo e massimo per gruppo
# -----------------------------
stats = (
    df.groupby(["OVERFLOW", "Elements"])["Latenza"]
    .agg(["mean", "min", "max"])
    .reset_index()
)

# -----------------------------
# Crea la figura
# -----------------------------
plt.figure(figsize=(10,6))
colors = plt.cm.tab10.colors

# -----------------------------
# Plot per ogni gruppo di 'Elements'
# -----------------------------
x_labels = sorted(stats["OVERFLOW"].unique())
n_groups = len(stats["Elements"].unique())
width = 0.15

for i, elements in enumerate(sorted(stats["Elements"].unique())):
    subset = stats[stats["Elements"] == elements]

    x_num = [x_labels.index(x_val) for x_val in subset["OVERFLOW"]]
    offset = (i - (n_groups-1)/2) * width
    x_dodge = [x + offset for x in x_num]

    y = subset["mean"]
    err_low = y - subset["min"]
    err_high = subset["max"] - y

    plt.errorbar(
        x_dodge,
        y,
        yerr=[err_low, err_high],
        fmt="o",
        label=f"Elements = {elements}",
        color=colors[i % len(colors)],
        capsize=4
    )

# -----------------------------
# Migliora la leggibilità
# -----------------------------
plt.xticks(range(len(x_labels)), x_labels, rotation=30, ha="right")
plt.xlabel("Overflow Strategy")
plt.ylabel("Average Latency [s]")
plt.title("Latency by Overflow Strategy")
plt.grid(True, axis="y")
plt.legend()
\end{lstlisting}

\section{Analisi dei dati}

L'analisi dei dati si concentra principalmente su due metriche fondamentali:
\begin{itemize}
    \item la latenza delle operazioni e dei task eseguiti;
    \item l'impatto sull'utilizzo della memoria.
\end{itemize}

L'obiettivo è valutare come le diverse configurazioni della hashtable influenzino le prestazioni complessive.

\subsection{Hash Data}

In questa sezione vengono analizzati i risultati relativi alle diverse configurazioni della hashtable.  
In particolare, si osserva l'influenza dei seguenti parametri:
\begin{itemize}
    \item strategia di HASHING;
    \item fattore di carico massimo (\textit{MAX LOAD FACTOR});
    \item strategia di gestione delle collisioni (\textit{OVERFLOW});
    \item politica di ridimensionamento (\textit{RESIZE}).
\end{itemize}

\subsubsection{Latenza}

\begin{figure}[H]
    \centering
    \includegraphics[width=0.8\textwidth]{hash_graphs/latency_by_overflow_strategy.pdf}
    \caption{Latenza in funzione della strategia di gestione delle collisioni (Integer e potenza di 10)}
\end{figure}

\begin{figure}[H]
    \centering
    \includegraphics[width=0.8\textwidth]{doubling/hash_graphs/latency_by_overflow_strategy.pdf}
    \caption{Latenza in funzione della strategia di gestione delle collisioni (Integer e doubling)}
\end{figure}

Dai grafici della latenza relativi alle strategie di OVERFLOW si evince che il tempo necessario a completare le operazioni rimane complessivamente 
simile al variare delle diverse tecniche.  
Un'eccezione è rappresentata dal metodo DOUBLE HASHING, che all'aumentare del numero di chiavi tende a peggiorare sia in termini di latenza media 
sia nei picchi massimi.

Tra i quattro metodi analizzati, CUCKOO e CHAINING mostrano inoltre una maggiore stabilità, con oscillazioni della latenza più contenute rispetto 
ai metodi di OPEN ADDRESSING.

Dal grafico generato dal doubling test inoltre si evince che la latenza circa raddoppia al raddoppiare della dimensione dell'input.

\begin{figure}[H]
    \centering
    \includegraphics[width=0.8\textwidth]{hash_graphs/latency_by_resize_strategy.pdf}
    \caption{Latenza in funzione della strategia di RESIZE (Integer e potenza di 10)}
\end{figure}

\begin{figure}[H]
    \centering
    \includegraphics[width=0.8\textwidth]{doubling/hash_graphs/latency_by_resize_strategy.pdf}
    \caption{Latenza in funzione della strategia di RESIZE (Integer e doubling)}
\end{figure}

Per quanto concerne la latenza in relazione alle strategie di RESIZE, si osserva un comportamento coerente con le aspettative teoriche.  
La strategia CONSTANT risulta la meno efficiente: all'aumentare del numero di elementi, e quindi allontanandosi dalla dimensione iniziale fissata, 
le prestazioni peggiorano sensibilmente.

Le strategie DOUBLE e DYNAMIC presentano invece prestazioni complessivamente simili, con una latenza media leggermente migliore nel caso della strategia DOUBLE.

Dal grafico generato dal doubling test inoltre si evince che la latenza circa raddoppia al raddoppiare della dimensione dell'input.

\begin{figure}[H]
    \centering
    \includegraphics[width=0.8\textwidth]{hash_graphs/latency_by_hash_strategy.pdf}
    \caption{Latenza in funzione della strategia di HASHING (Integer e potenza di 10)}
\end{figure}

\begin{figure}[H]
    \centering
    \includegraphics[width=0.8\textwidth]{doubling/hash_graphs/latency_by_hash_strategy.pdf}
    \caption{Latenza in funzione della strategia di HASHING (Integer e doubling)}
\end{figure}

Per quanto riguarda le due strategie di HASHING considerate, si osserva una differenza marcata: la prima risulta nettamente più efficiente 
in termini di latenza rispetto alla seconda.

\begin{figure}[H]
    \centering
    \includegraphics[width=0.8\textwidth]{hash_graphs/latency_by_max_load_factor_strategy.pdf}
    \caption{Latenza in funzione del MAX LOAD FACTOR (Integer e potenza di 10)}
\end{figure}

\begin{figure}[H]
    \centering
    \includegraphics[width=0.8\textwidth]{doubling/hash_graphs/latency_by_max_load_factor_strategy.pdf}
    \caption{Latenza in funzione del MAX LOAD FACTOR (Integer e doubling)}
\end{figure}

Il MAX LOAD FACTOR sembra influire in maniera limitata sulla latenza complessiva: le differenze osservate tra le varie configurazioni risultano contenute.

\subsubsection{Memoria}

\begin{figure}[H]
    \centering
    \includegraphics[width=0.8\textwidth]{hash_graphs/memory_by_overflow_strategy.pdf}
    \caption{Utilizzo della memoria in funzione della strategia di gestione delle collisioni (Integer e potenza di 10)}
\end{figure}

\begin{figure}[H]
    \centering
    \includegraphics[width=0.8\textwidth]{doubling/hash_graphs/memory_by_overflow_strategy.pdf}
    \caption{Utilizzo della memoria in funzione della strategia di gestione delle collisioni (Integer e doubling)}
\end{figure}

A differenza della latenza, l'utilizzo della memoria varia in modo significativo tra le diverse strategie.  
I metodi CHAINING e CUCKOO richiedono una quantità di memoria sensibilmente maggiore rispetto alle tecniche di OPEN ADDRESSING, 
sia in media sia nei picchi massimi delle singole esecuzioni.

\begin{figure}[H]
    \centering
    \includegraphics[width=0.8\textwidth]{hash_graphs/memory_by_resize_strategy.pdf}
    \caption{Utilizzo della memoria in funzione della strategia di RESIZE (Integer e potenza di 10)}
\end{figure}

\begin{figure}[H]
    \centering
    \includegraphics[width=0.8\textwidth]{doubling/hash_graphs/memory_by_resize_strategy.pdf}
    \caption{Utilizzo della memoria in funzione della strategia di RESIZE (Integer e doubling)}
\end{figure}

Per quanto riguarda le strategie di RESIZE, si osserva che il metodo CONSTANT tende a occupare più memoria rispetto agli altri due approcci.  
Tra DOUBLE e DYNAMIC, il primo mostra in media un comportamento leggermente più efficiente.

\begin{figure}[H]
    \centering
    \includegraphics[width=0.8\textwidth]{hash_graphs/memory_by_hash_strategy.pdf}
    \caption{Utilizzo della memoria in funzione della strategia di HASHING (Integer e potenza di 10)}
\end{figure}

\begin{figure}[H]
    \centering
    \includegraphics[width=0.8\textwidth]{doubling/hash_graphs/memory_by_hash_strategy.pdf}
    \caption{Utilizzo della memoria in funzione della strategia di HASHING (Integer e doubling)}
\end{figure}

La strategia di HASHING influisce in modo rilevante sull'utilizzo della memoria.  
In particolare, si osserva un comportamento opposto rispetto alla latenza: la seconda strategia risulta più efficiente della prima 
dal punto di vista dell'occupazione di memoria.

\begin{figure}[H]
    \centering
    \includegraphics[width=0.8\textwidth]{hash_graphs/memory_by_max_load_factor_strategy.pdf}
    \caption{Utilizzo della memoria in funzione del MAX LOAD FACTOR (Integer e potenza di 10)}
\end{figure}

\begin{figure}[H]
    \centering
    \includegraphics[width=0.8\textwidth]{doubling/hash_graphs/memory_by_max_load_factor_strategy.pdf}
    \caption{Utilizzo della memoria in funzione del MAX LOAD FACTOR (Integer e doubling)}
\end{figure}

Il MAX LOAD FACTOR, come atteso, incide in modo diretto sull'utilizzo della memoria: valori più bassi comportano un maggiore consumo di spazio, 
mentre configurazioni standard o elevate risultano più efficienti.

\paragraph{Sintesi dei risultati}

Dall'analisi complessiva emerge come le diverse configurazioni della hashtable influenzino in maniera differente le metriche considerate, 
evidenziando un compromesso tra latenza e utilizzo della memoria.

Per quanto riguarda la latenza, le strategie di OVERFLOW mostrano comportamenti generalmente simili, ad eccezione del DOUBLE HASHING, 
che tende a degradare all'aumentare del numero di chiavi. Le tecniche CUCKOO e CHAINING garantiscono invece una maggiore stabilità.  
Tra le politiche di RESIZE, la strategia CONSTANT risulta inefficiente, mentre DOUBLE e DYNAMIC offrono prestazioni comparabili, 
con un leggero vantaggio per DOUBLE. La scelta della funzione di HASHING risulta determinante, mentre il MAX LOAD FACTOR ha un impatto limitato.
Generalmente inoltre si evince un andamento della latenza che tende a raddoppiare al raddoppio della dimensione dell'input.

Dal punto di vista della memoria, emergono differenze più marcate: CHAINING e CUCKOO risultano più costosi rispetto all'OPEN ADDRESSING, 
mentre la strategia CONSTANT comporta un maggiore utilizzo di memoria rispetto alle alternative. Inoltre, la seconda strategia di HASHING 
si dimostra più efficiente, evidenziando un comportamento opposto rispetto alla latenza.

In sintesi, non esiste una configurazione universalmente ottimale: la scelta dipende dal contesto applicativo e dai vincoli specifici. 
Tuttavia, l'utilizzo della strategia DOUBLE per il RESIZE e una scelta accurata della funzione di HASHING rappresentano, in generale, 
un buon compromesso tra efficienza e stabilità.

\subsection{Operation Data}

Questa sezione è dedicata all'analisi delle singole operazioni sulla hashtable (inserimento, ricerca, aggiornamento e cancellazione).  
Viene studiato il comportamento delle operazioni al variare delle diverse configurazioni, con particolare attenzione a:
\begin{itemize}
    \item distribuzione delle chiavi;
    \item strategia di hashing;
    \item fattore di carico massimo (\textit{max load factor});
    \item strategia di gestione delle collisioni (\textit{overflow});
    \item politica di ridimensionamento (\textit{resize}).
\end{itemize}

\subsubsection{Latenza}

\begin{figure}[H]
    \centering
    \includegraphics[width=0.8\textwidth]{op_graphs/latency_by_operations_&_overflow_strategy.pdf}
    \caption{Latenza per operazione in funzione della strategia di gestione delle collisioni (Integer e potenza di 10)}
\end{figure}

\begin{figure}[H]
    \centering
    \includegraphics[width=0.8\textwidth]{doubling/op_graphs/latency_by_operations_&_overflow_strategy.pdf}
    \caption{Latenza per operazione in funzione della strategia di gestione delle collisioni (Integer e doubling)}
\end{figure}

Dal grafico si osserva come l'operazione di \textit{insert} presenti generalmente una latenza superiore rispetto alle altre. 
Ciò è dovuto al controllo del \textit{load factor} e alla possibile necessità di effettuare operazioni di \textit{resize}, che introducono un costo aggiuntivo.

Tra le diverse strategie di gestione delle collisioni, il \textit{double hashing} mostra prestazioni significativamente peggiori 
su tutte le operazioni, in particolare sull'inserimento, dove il costo delle sonde multiple incide maggiormente.

Le strategie \textit{cuckoo hashing} e \textit{chaining} evidenziano prestazioni complessivamente simili. Tuttavia, si nota che:
\begin{itemize}
    \item il \textit{cuckoo hashing} presenta picchi leggermente più contenuti per l'operazione di insert all'aumentare del numero di chiavi;
    \item il \textit{chaining} risulta leggermente più stabile nelle operazioni di ricerca, aggiornamento e cancellazione.
\end{itemize}

L'\textit{open addressing} con probing lineare, pur mostrando una latenza media comparabile alle precedenti strategie, 
presenta prestazioni peggiori per dataset di piccole dimensioni. Tuttavia, al crescere del numero di chiavi, mostra una maggiore stabilità 
e una crescita più controllata della latenza.

\begin{figure}[H]
    \centering
    \includegraphics[width=0.8\textwidth]{op_graphs/latency_by_operations_&_resize_strategy.pdf}
    \caption{Latenza per operazione in funzione della strategia di ridimensionamento (Integer e potenza di 10)}
\end{figure}

\begin{figure}[H]
    \centering
    \includegraphics[width=0.8\textwidth]{doubling/op_graphs/latency_by_operations_&_resize_strategy.pdf}
    \caption{Latenza per operazione in funzione della strategia di ridimensionamento (Integer e doubling)}
\end{figure}

L'analisi delle politiche di \textit{resize} evidenzia come la strategia \textit{constant} comporti latenze medie e picchi significativamente 
superiori rispetto alle altre. Questo è dovuto alla maggiore frequenza di ridimensionamenti o alla loro scarsa efficienza.

Le strategie \textit{double} e \textit{dynamic} mostrano invece prestazioni medie simili. Tuttavia:
\begin{itemize}
    \item la strategia \textit{dynamic} garantisce una maggiore stabilità complessiva;
    \item la strategia \textit{double} presenta picchi più marcati, soprattutto all'aumentare del numero di chiavi.
\end{itemize}

\begin{figure}[H]
    \centering
    \includegraphics[width=0.8\textwidth]{op_graphs/latency_by_operation_hash_strategy.pdf}
    \caption{Latenza per operazione in funzione della funzione di hash (Integer e potenza di 10)}
\end{figure}

\begin{figure}[H]
    \centering
    \includegraphics[width=0.8\textwidth]{doubling/op_graphs/latency_by_operation_hash_strategy.pdf}
    \caption{Latenza per operazione in funzione della funzione di hash (Integer e doubling)}
\end{figure}

Il confronto tra le due funzioni di hash mostra che la prima strategia risulta più efficiente sia in termini di latenza media sia nei picchi massimi, 
soprattutto su dataset di grandi dimensioni.

La seconda funzione di hash, pur essendo meno performante in media, mostra risultati leggermente migliori nel caso di dataset molto piccoli, 
dove il costo computazionale della funzione ha un impatto più rilevante rispetto alla distribuzione delle chiavi.

\begin{figure}[H]
    \centering
    \includegraphics[width=0.8\textwidth]{op_graphs/latency_by_operations_&_max_load_factor.pdf}
    \caption{Latenza per operazione in funzione del max load factor (Integer e potenza di 10)}
\end{figure}

\begin{figure}[H]
    \centering
    \includegraphics[width=0.8\textwidth]{doubling/op_graphs/latency_by_operations_&_max_load_factor.pdf}
    \caption{Latenza per operazione in funzione del max load factor (Integer e doubling)}
\end{figure}

Le prestazioni in funzione del \textit{max load factor} risultano nel complesso simili tra le diverse configurazioni. Tuttavia, emergono alcune differenze:
\begin{itemize}
    \item un \textit{load factor} basso (\textit{LOW}) garantisce maggiore stabilità nelle varie operazioni;
    \item valori più elevati (\textit{STANDARD} e \textit{HIGH}) offrono migliori prestazioni medie nell'operazione di insert, 
    grazie a un minor numero di ridimensionamenti.
\end{itemize}

Questo evidenzia il classico trade-off tra utilizzo della memoria e prestazioni.

\begin{figure}[H]
    \centering
    \includegraphics[width=0.8\textwidth]{op_graphs/latency_insert_exist_by_load_factor_&_all_overflows.pdf}
    \caption{Latenza delle'operazioni di insert e get al variare del load factor e della strategia di overflow (Integer e potenza di 10)}
\end{figure}

\begin{figure}[H]
    \centering
    \includegraphics[width=0.8\textwidth]{doubling/op_graphs/latency_insert_exist_by_load_factor_&_all_overflows.pdf}
    \caption{Latenza delle'operazioni di insert e get al variare del load factor e della strategia di overflow (Integer e doubling)}
\end{figure}

Le strategie \textit{chaining} e \textit{cuckoo hashing} mantengono latenze inferiori e più stabili, mentre il \textit{double hashing} 
continua a mostrare prestazioni peggiori.

In generale, l'aumento del \textit{load factor} tende a penalizzare maggiormente le tecniche basate su \textit{open addressing}, 
a causa dell'incremento delle collisioni e del numero medio di sonde necessarie.

\paragraph{Sintesi dei risultati}

L'analisi della latenza delle operazioni evidenzia come le prestazioni di una hashtable dipendano in modo significativo dall'interazione 
tra strategia di gestione delle collisioni, politica di ridimensionamento e valore del \textit{max load factor}.

In generale, le operazioni di inserimento risultano le più costose, principalmente a causa del controllo del fattore di carico e delle eventuali 
operazioni di \textit{resize}. Tra le strategie di \textit{overflow}, il \textit{double hashing} si dimostra la meno efficiente, mostrando latenze 
elevate e scarsa scalabilità. Al contrario, \textit{chaining} e \textit{cuckoo hashing} offrono un buon compromesso tra prestazioni medie e stabilità, 
mentre il probing lineare risulta più sensibile alla dimensione del dataset ma garantisce una crescita della latenza più regolare.

Per quanto riguarda il \textit{resize}, strategie adattive come la \textit{dynamic} permettono di contenere i picchi di latenza, risultando 
complessivamente più robuste rispetto a strategie più semplici come la \textit{constant}. Anche la scelta della funzione di hash incide in modo rilevante: 
una buona distribuzione delle chiavi diventa sempre più critica all'aumentare della dimensione dei dati.

Infine, il \textit{max load factor} introduce un chiaro compromesso tra efficienza temporale e utilizzo della memoria: valori bassi migliorano 
la stabilità e riducono le collisioni, mentre valori più elevati riducono la frequenza dei ridimensionamenti ma possono degradare le prestazioni 
nelle tecniche di \textit{open addressing}.

Complessivamente, i risultati suggeriscono che non esiste una configurazione ottimale in senso assoluto, ma che la scelta dei parametri debba essere 
guidata dal contesto applicativo, dal volume dei dati e dai requisiti di prestazione.


\subsection{Mixed Data}

In questa sezione vengono analizzati cinque grafici complessivi che confrontano direttamente le prestazioni delle diverse tipologie di dati.

L'obiettivo è valutare il comportamento della hashtable al variare del tipo di input, evidenziando eventuali differenze in termini di 
latenza e utilizzo della memoria.

\begin{figure}[H]
    \centering
    \includegraphics[width=0.8\textwidth]{mix_graphs/latency_by_keytype_overflow.pdf}
    \caption{Latenza per operazione in funzione della strategia di gestione delle collisioni, suddivisa per tipo di chiave. (Integer e potenza di 10)}
\end{figure}

\begin{figure}[H]
    \centering
    \includegraphics[width=0.8\textwidth]{doubling/mix_graphs/latency_by_keytype_overflow.pdf}
    \caption{Latenza per operazione in funzione della strategia di gestione delle collisioni, suddivisa per tipo di chiave. (Integer e doubling)}
\end{figure}

Dal grafico emerge chiaramente che le chiavi di tipo \textit{Integer} garantiscono le migliori prestazioni complessive.  
Le \textit{String} generate casualmente presentano latenze leggermente superiori, ma la differenza rispetto agli \textit{Integer} è minima.  
Al contrario, le chiavi basate su \textit{Real Words} mostrano latenze significativamente più elevate, probabilmente a causa di una distribuzione 
meno uniforme che genera più collisioni.

\begin{figure}[H]
    \centering
    \includegraphics[width=0.8\textwidth]{mix_graphs/latency_by_keytype_resize.pdf}
    \caption{Latenza per operazione in funzione della strategia di \textit{resize}, suddivisa per tipo di chiave. (Integer e potenza di 10)}
\end{figure}

\begin{figure}[H]
    \centering
    \includegraphics[width=0.8\textwidth]{doubling/mix_graphs/latency_by_keytype_resize.pdf}
    \caption{Latenza per operazione in funzione della strategia di \textit{resize}, suddivisa per tipo di chiave. (Integer e doubling)}
\end{figure}

Risultati analoghi si osservano analizzando la strategia di \textit{resize}.  
Gli \textit{Integer} mantengono prestazioni superiori, seguiti dalle \textit{String} casuali. Le chiavi basate su \textit{Real Words} mostrano 
latenze sensibilmente più alte, confermando il maggiore impatto delle collisioni su dataset meno uniformi.

\begin{figure}[H]
    \centering
    \includegraphics[width=0.8\textwidth]{mix_graphs/latency_by_keytype_hash.pdf}
    \caption{Latenza per operazione in funzione della funzione di hash, suddivisa per tipo di chiave. (Integer e potenza di 10)}
\end{figure}

\begin{figure}[H]
    \centering
    \includegraphics[width=0.8\textwidth]{doubling/mix_graphs/latency_by_keytype_hash.pdf}
    \caption{Latenza per operazione in funzione della funzione di hash, suddivisa per tipo di chiave. (Integer e doubling)}
\end{figure}

Anche la scelta della funzione di hash evidenzia la stessa tendenza: le chiavi \textit{Integer} ottimizzano la distribuzione, riducendo le collisioni 
e quindi la latenza, mentre le \textit{Real Words} risultano più penalizzate.

\begin{figure}[H]
    \centering
    \includegraphics[width=0.8\textwidth]{mix_graphs/latency_by_operation_keytype_elements.pdf}
    \caption{Latenza per singola operazione in funzione del tipo di chiave e del numero di elementi. (Integer e potenza di 10)}
\end{figure}

\begin{figure}[H]
    \centering
    \includegraphics[width=0.8\textwidth]{doubling/mix_graphs/latency_by_operation_keytype_elements.pdf}
    \caption{Latenza per singola operazione in funzione del tipo di chiave e del numero di elementi. (Integer e doubling)}
\end{figure}

A livello delle singole operazioni, le prestazioni confermano quanto osservato precedentemente.  
Si nota un leggero vantaggio dell'approccio ibrido (\textit{Hybrid}) nell'operazione di \textit{insert}, rispetto alle strategie full stringa, 
grazie a una combinazione più efficiente di tipi di chiave.

\begin{figure}[H]
    \centering
    \includegraphics[width=0.8\textwidth]{mix_graphs/memory_by_key_type.pdf}
    \caption{Memoria media utilizzata in funzione del tipo di chiave. (Integer e potenza di 10)}
\end{figure}

\begin{figure}[H]
    \centering
    \includegraphics[width=0.8\textwidth]{doubling/mix_graphs/memory_by_key_type.pdf}
    \caption{Memoria media utilizzata in funzione del tipo di chiave. (Integer e doubling)}
\end{figure}

L'analisi della memoria evidenzia che gli \textit{Integer} richiedono meno spazio, seguiti dall'approccio ibrido, mentre le \textit{String} 
complete risultano le più costose in termini di consumo di memoria.

\paragraph{Sintesi dei risultati}

Dall'analisi dei grafici emerge chiaramente come il tipo di chiave influisca in modo determinante sulle prestazioni della hashtable.  
Gli \textit{Integer} offrono i migliori risultati sia in termini di latenza sia di memoria utilizzata, grazie alla loro natura semplice e 
alla distribuzione uniforme.  
Le \textit{String} casuali mantengono prestazioni simili ma leggermente peggiori, mentre le chiavi basate su \textit{Real Words} 
risultano più penalizzate a causa di collisioni più frequenti e maggiore overhead di memoria.

L'approccio ibrido rappresenta un buon compromesso, combinando la compattezza degli \textit{Integer} con la flessibilità delle \textit{String}, 
offrendo un leggero miglioramento nelle operazioni di inserimento rispetto alle strategie full stringa.

In generale, i risultati suggeriscono che la scelta del tipo di chiave deve essere guidata dal contesto applicativo e dai requisiti di prestazione, 
privilegiando tipologie più compatte e uniformi quando si cerca massima efficienza.